%%% Uso da classe LaTeX Ifes8 para elaboração de trabalho acadêmicos
%%% de acordo como as normas estabelecidas no documento “Normas para
%%% Apresentação de Trabalhos Acadêmicos e Científicos: Documento Impresso e/ou
%%% Digital”, 9ª edição -- 2024.
%%% 
%%% Para a versão mais recente desta classe verifique:
%%% https://github.com/quigonjeff/ifes8
%%%
%%% Autor: Prof. Jefferson O. Andrade <jefferson.andrade@ifes.edu.br>
%%% Atualizado por: Gustavo Cardoso Pimentel em 11 de agosto de 2026 (Adequação às normas Ifes 2024 - 9ª edição)
%%% Data de criação: 2021-07-01 | Data de atualização: 2026-08-11
%%%
\documentclass[times,english,brazilian,oneside,section=TITLE]{ifes8}

\usepackage[utf8]{inputenc}
\usepackage{lastpage}           % Usado pelo exemplo de ficha catalográfica
\usepackage[alf,abnt-full-initials=yes]{abntex2cite} % Esta opção é exigida pelas normas do ifes.

% Ajuste ABNT NBR 10520:2023 (Citações autor-data em caixa mista)
\makeatletter
\renewcommand{\citeauthor}{\citeauthoronline}
\makeatother

\usepackage{microtype}          % para melhorias de justificação
\usepackage{morefloats}         % permite mais floats
\usepackage{blindtext}          % para gerar texto aleatório (dummy text)
\usepackage{tikz}
\usetikzlibrary[topaths]
\usepackage{mathtools}
\usepackage{xcolor}
\usepackage{algorithm}
\usepackage{algpseudocode}
\usepackage{textcomp}
\usepackage{eurosym}
\usepackage{listings}
\input{configlistings}

%%% Definição da linguagem padrão do documento (pacote babel) para
%%% definir que certas porções do texto (como o resumo, por exemplo)
%%% estão em uma língua estrangeira, usar as macros
%%% \foreignlanguage{languageB}{Text in another language}
%%% ou
%%% \begin{otherlanguage}{languageB}
%%% ...
%%% \end{otherlanguage}
\selectlanguage{brazilian}

%%% Ajuste do alinhamento do caption de figuras, tabelas e quadros.
% \captionstyle{\raggedright}
% \normalcaption
% \normalcaptionwidth
% \captionwidth{\linewidth}


%%% Define que todos os códigos fontes construídos com o ambiente
%%% `lstlisting' terão uma borda simples.
\lstset{frame=single}

\titulo{Metodologia para Caracterização Viscoelástica e Dimensionamento de Enrijecedores à Flexão para Risers Flexíveis}

\autor{Alexandre Lima Simões\\Pedro Antônio Gomes Goronci}
\autorficha{Simões, Alexandre Lima; Goronci, Pedro Antônio Gomes}

\orientador{Prof. Dr. Vinicius Erler de Souza Ramos}

% \coorientadora{Prof.ª Dr.ª Fulana Detal Longname}

\instituicao{Instituto Federal do Espírito Santo}

\curso{Bacharelado em Engenharia Mecânica}

\data{2026}

\local{Vitória, ES}

\preambulo{Trabalho de Conclusão de Curso apresentado à Coordenadoria
  do Curso de Engenharia Mecânica do Instituto Federal do Espírito
  Santo, Campus Serra, como requisito parcial para a obtenção do
  título de Bacharel em Engenharia Mecânica.}

\tipotrabalho{Trabalho de Conclusão de Curso}



\begin{document}

%%% A macro \pretextual não precisa ser explicitamente chamada porque
%%% o abntex2 aciona essa macro no início do documento.
% \pretextual

\imprimircapa
\addtocounter{page}{-1}

\imprimirfolhaderosto*

\begin{fichacatalografica}
  \vspace*{\fill}
  % Posição vertical
  \begin{center}
    % Minipage Centralizado
    \fbox{
      \begin{minipage}[t]{12.5cm}
        \vspace*{3mm}
        \begin{minipage}[t]{2cm}
          X000x
        \end{minipage}
        \begin{minipage}[t]{10cm}
          \imprimirautorficha.\vspace{2mm}

          \hspace{0.5cm}\imprimirtitulo\ / \imprimirautor. ---
          \imprimirlocal, \imprimirdata.\\ 

          \hspace{0.5cm}\pageref{LastPage} p.:\ il. (algumas color.);
          30 cm.\\ 

          \hspace{0.5cm}\imprimirorientadorRotulo~\imprimirorientador.\\

          \hspace{0.5cm}\imprimirtipotrabalho\ ---
          \imprimirinstituicao, \imprimirdata\\

          \hspace{0.5cm}        
          1. Viscoelasticidade.
          2. Engenharia do petróleo.
          3. Equações diferenciais.
          4. Métodos numéricos.
          I. Vinicius Erler de Souza Ramos.
          II. Instituto Federal do Espírito Santo.
          III. Campus Vitória.
          IV. Título.
          
          \begin{flushright}
            CDU 00:000:000.0
          \end{flushright}
          \vspace*{1mm}
        \end{minipage}
      \end{minipage}
    }
    
    \vspace*{3mm}
    {\footnotesize Bibliotecário(a) responsável: Nome Completo -- CRB X/XXXX}
  \end{center}
\end{fichacatalografica}
\addtocounter{page}{-1}

%%% ====================================================================
%%% Exemplo de construção da folha de aprovação. Depois da
%%% apresentação do trabalho, quando a folha de apresentação real
%%% tiver sido assinada pela banca, a folha real deve ser digitalizada
%%% e o ambiente abaixo deve ser substituído por:
%%% \includepdf{folhadeaprovacao_final.pdf}
\begin{folhadeaprovacao}
  \begin{center}
    {\ABNTEXchapterfont\MakeTextUppercase{\imprimirautor}}\\
    \begin{center}
      \ABNTEXchapterfont\MakeTextUppercase{\imprimirtitulo}\\
    \end{center}
    \hspace{.45\textwidth}
    \begin{minipage}{.5\textwidth}
      {\footnotesize{\imprimirpreambulo}}\\
    \end{minipage}%
  \end{center}
  % \vspace{-1cm}%
  \begin{center}
    Aprovado em 11 de agosto de 2026.\\[15mm]
    \textbf{COMISSÃO EXAMINADORA}\\[5mm]
    \assinatura{\imprimirorientador \\
      Instituto Federal do Espírito Santo \\
      Orientador}
    % \vspace{-1cm}
    \makeatletter
    \abntex@ifnotempty{\imprimircoorientador}{%
      \assinatura{\imprimircoorientador \\
        Instituto Federal do Espírito Santo \\
        Coorientador}
    }
    \makeatother
    % \vspace{-1cm}
    \assinatura{Membro Convidado da Banca 1 \\
      Instituição Externa 1}
    % \vspace{-1cm}
    \assinatura{Membro Convidado da Banca 2 \\
      Instituição Externa 2}
  \end{center}
\end{folhadeaprovacao}


%%% --------------------------------------------------------------------
%%% Ambiente para escrita da declaração do autor
%%% 2020-09-02: Jefferson Andrade. A declaração do autor não é mais necessária
%%% nos trabalhos acadêmcios do Ifes.
% \begin{declaracaodoautor}

%   \vspace*{1.5cm}

%   Declaro, para fins de pesquisa acadêmica, didática e
%   técnico-científica, que este Trabalho de Conclusão de Curso pode ser
%   parcialmente utilizado, desde que se faça referência à fonte e ao
%   autor.

%   \vspace*{2.5cm}

%   \centering

%   \imprimirlocal, 11 de agosto de 2026.

%   \vspace*{2.5cm}

%   \imprimirautor

%   \vspace*{\fill}
  
% \end{declaracaodoautor}


%%% ====================================================================
%%% Ambiente para a escrita da dedicatória.
\begin{dedicatoria}
  \vspace*{\fill}
  \hspace{0.2\textwidth}
  \begin{minipage}{0.8\textwidth}
    Este trabalho é dedicado aos que acreditam que Pikachu sempre será o
    melhor pokemon.
  \end{minipage}
\end{dedicatoria}


%%% ====================================================================
%%% Agradecimentos
\begin{agradecimentos}
  Agradeço a todos os que fizeram alguma coisa pela qual eu deveria
  ser grato.
\end{agradecimentos}


%%% ====================================================================
%%% Epígrafe
\begin{epigrafe}
  \vspace*{\fill}
  \begin{otherlanguage}{english}
    \begin{flushright}
      \begin{SingleSpace}
        People who think they know everything are a great \\
        annoyance to those of us who do. \cite{Asimov2004}
      \end{SingleSpace}
    \end{flushright}
  \end{otherlanguage}
\end{epigrafe}


%%% Ambiente para resumo em português
\begin{resumo}

  A engenharia \emph{offshore} (fora da costa) é responsável pelo desenvolvimento de atividades
  relacionadas à exploração e produção de petróleo e gás natural em ambientes marítimos. Para essas
  aplicações, o uso de linhas flexíveis é essencial para o escoamento da produção e para a interligação
  de equipamentos marinhos com a plataforma de operação FPSO (\emph{Floating Production Storage and Offloading}).
  De modo geral, as linhas flexíveis não possuem resistência à flexão suficiente para suportar as cargas associadas
  ao ambiente que são trasnferidas para as conexões de topo, sendo necessário o uso de enrijecedores à flexão.
  O objetivo deste trabalho é apresentar uma metodologia para a caracterização viscoelástica e 
  o dimensionamento de enrijecedores à flexão para risers flexíveis por meio da implementação 
  de um modelo matemático capaz de descrever o comportamento viscoelástico mediante regressão
  dos dados experimentais obtidos no ensaio de relaxação do poliuretano e implementação de solução
  numérica para avaliação da resposta dinâmica do sistema enrijecedor de curvatura com duto flexível
  em regime harmônico.

  Palavras-chave: Engenharia offshore. Enrijecedores à Flexão. Viscoelasticidade. Métodos numéricos. Dimensionamento. 
\end{resumo}

%%% Ambiente para resumo em inglês.
\begin{resumo}[Abstract]
  \begin{otherlanguage}{english}
    The offshore engineering is responsible for the development of activities related to the exploration and 
    production of oil and natural gas in marine environments. For these applications, the use of flexible 
    risers is essential for the flow of production and for the interconnection of marine equipment with the 
    FPSO (\emph{Floating Production Storage and Offloading}) platform. In general, flexible risers do not 
    have sufficient bending resistance to withstand the loads associated with the environment, which are transferred 
    to the top connections, making the use of bending stiffeners necessary. The objective of this work is to 
    present a methodology for the viscoelastic characterization and the design of bending stiffeners for flexible 
    risers through the implementation of a mathematical model capable of describing the viscoelastic behavior by 
    regression of the experimental data obtained in the polyurethane relaxation test and the implementation of a 
    numerical solution for evaluating the dynamic response of the bending stiffener system with flexible pipe under 
    harmonic regime.

    Keywords: Offshore Engineering. Bending Stiffeners. Viscoelasticity. Numerical Methods. Design.
  \end{otherlanguage}
\end{resumo}


% %%% Lista de figuras
% \renewcommand{\listfigurename}{Lista de figuras}
% \pdfbookmark[0]{\listfigurename}{lof}
% \listoffigures*
% \cleardoublepage


% %%% Lista de tabelas
% \pdfbookmark[0]{\listtablename}{lot}
% \listoftables*
% \cleardoublepage


% %%% Lista de quadros
% \pdfbookmark[0]{\listadequadrosname}{loq}
% \listadequadros*
% \cleardoublepage


% %%% Lista de abreviaturas
% \begin{abreviaturas}
%   \simb{Abrev.} -- Abreviatura
%   \simb{Assoc.} -- Associação
%   \simb{Atm.} -- Atmosfera
%   \simb{Bel.} -- Bacharel
%   \simb{Bioq.} -- Bioquímica
%   \simb{Cit.} -- Citação
%   \simb{Compl.} -- Complemento
%   \simb{Dic.} -- Dicionário
%   \simb{Dipl.} -- Diploma
% \end{abreviaturas}


% %%% Lista de siglas
% \begin{siglas}
%   \simb{ABNT} -- Associação Brasileira de Normas Técnicas
%   \simb{ACM} -- Association for Computing Machinery
%   \simb{APA} -- American Psychological Association
%   \simb{BIOS} -- Basic Input / Output System
%   \simb{CREA} -- Conselho Regional de Engenharia e Agronomia
%   \simb{IBGE} -- Instituto Brasileiro de Geografia e Estatística
%   \simb{IEEE} -- Institute of Electrical and Electronics Engineers
%   \simb{Ifes} -- Instituo Federal de Educação, Ciência e Tecnologia do Espírito Santo
%   \simb{ISDN} -- Integrated Services Digital Network
%   \simb{RISC} -- Reduced Instruction Set Computer
% \end{siglas}


%%% Lista de símbolos
\begin{simbolos}
  \simb{$\Gamma$} -- Letra grega Gama
  \simb{$\Lambda$} -- Lambda
  \simb{$\zeta$} -- Letra grega minúscula zeta
  \simb{$\in$} -- Pertence
  \simb{$\top$} -- Valor lógico máximo dentro de um reticulado regular booleano ou quasi-booleano.
\end{simbolos}
\cleardoublepage


%%% Sumário --- Table of Contents
\imprimirsumario
\cleardoublepage


%%% ====================================================================
%%% Início da parte textual do documento.
\textual


%%% Configuração do espaçamento entre títulos e texto
\setlength{\afterchapskip}{0\baselineskip}

% ----------------------------------------------------------------------
% Pedro
\chapter{Introdução}
\label{cha:introducao}

% Introduz o contexto.
A produção de petróleo e gás natural em alto-mar faz parte de cadeias ligadas ao fornecimento de energia e 
matérias-primas utilizadas em diversos setores da economia mundial. Para que essa produção ocorra com
segurança e confiabilidade, o projeto de sistemas e equipamentos submarinos devem atender exigências 
relacionadas às condições extremas de instalação e operação.

% Relaciona com linhas flexíveis e explica os diferentes trechos.
Nesse contexto, a comunicação entre as plataformas que operam longe da costa (\emph{offshore}) e os sistemas
submarinos que realizam as operações no leito marinho é um objeto de estudo relevante para a engenharia. As 
unidades flutuantes de operação e armazenamento (\emph{FPSO -- Floating Production Storage and Offloading})
são conectadas por meio de linhas flexíveis que permitem o transporte de petróleo e gás natural diretamente 
da extração por meio de dutos flexíveis, e existem também cabos desenvolvidos para a transmissão de sinais, 
dados e energia entre as plataformas e os sistemas submarinos, denominados cabos umbilicais. O trecho suspenso
que em mar que vai do ponto de conexão da plataforma até a zona de contato com o leito marinho é comumente 
atribuído pelo nome de \emph{riser} e é caracterizado pela forte resposta dinâmica, enquanto o trecho estático
é o restante da linha flexível que permanece em contato com o leito marinho e é chamado de \emph{flowline}.

% Destaca a severidade do ambiente offshore.
O trecho dinâmico das linhas flexíveis está sempre sujeito à ação de forças hidrodinâmicas relacionadas à 
corrente marítima e à incidência de ondas, cargas provenientes da movimentação da plataforma e que podem ser
agravadas por condições ambientais adversas e também por variações na operação da plataforma, como cargas 
acidentais (e.g. perda de linhas de ancoragem, inundação de compartimentos da plataforma). Dessa forma, 
apesar das linhas flexíveis serem projetadas para suportar condições de elevada tração e curvatura, o projeto
desse tipo de produto não funcionaria sem o uso de acessórios específicos que possam contribuir para sua 
operação, considerando a severidade do ambiente \emph{offshore}. 

% Relaciona linhas flexíveis com a necessidade de enrijecedores à flexão.
De modo geral, na região de conexão do umbilical com a plataforma (região de topo), a flexão sempre é elevada,
podendo ultrapassar o valor de máxima curvatura permitida no projeto definido pela análise local na linha flexível.
Para evitar que o limite de curvatura seja excedido, utilizam-se de um equipamento capaz de suavizar os ângulos
gerados na região de topo a partir do aumento de rigidez com uma redução gradual ao longo do comprimento
da linha. Esses acessórios são chamados de enrijecedores à flexão e possuem um formato cônico com a base maior 
próxima do ponto de conexão e envolvem o cabo submarino, de modo que o atravessa o centro do enrijecedor. 

% Importancia do dimensionamento do enrijecedor.
Dada a função de suavização da curvatura na região de topo, o enrijecedor à flexão desempenha um papel crucial 
para a segurança estrutural da linha flexível, sendo bastante presente em projetos que envolvam cabos submarinos.
Ademais, por se tratar de um componente que dependendo das condições de operação pode ser bem exigente e específico,
seu custo é bem elevado e pode representar um custo determinante para a execução ou não de um projeto. Portanto, 
o dimensionamento correto dos enrijecedores é apreciado tanto do ponto de vista da engenharia, por conta de sua 
relevância estrutural, quanto do ponto de vista econômico, devido a sua relevância e criticidade.

% Características do enrijecedor.
Justificada a importância dos enrijecedores, ainda se faz necessário compreender como suas propriedades mecânicas 
devem atender os requisitos funcionais do sistema enrijecedor e linha flexível. São feitos de poliuretano -- um 
material polimérico de elevada resistência mecânica, resistência abrasiva e capacidade de absorvação de impactos 
\cite{Balmoral?} --, possuem um inserto metálico no seu interior para permitir a união com outros acessórios importantes do 
sistema de fixação da linha flexível e são usinados para adquirir a forma cônica adequada.

% A importância de simular o comportamento estrutural do enrijecedor à flexão.

% Problemas de uma abordagem elástica.

% Vantagens de uma abordagem viscoelástica.

% Necessidade de métodos numéricos

% Noções importantes (que vão ser abordadas para justificar as considerações)

% \begin{itemize}
%   \item Contextualização da indústria offshore e a importância das linhas flexíveis.
%   \item Problema de pesquisa: necessidade de enrijecedores à flexão para risers flexíveis.
%   \item Aprofundamento do problema: desafios na caracterização viscoelástica de poliuretanos e dimensionamento de enrijecedores.
%   \item Análise global e inputs do problema (cargas submetidas).
%   \item Apresentação de ferramentas e métodos numéricos relevantes para o desenvolvimento da análise estrutural do \emph{bend stiffener}.
%   \begin{itemize}
%     \item Softwares de análise estrutural e simulação numérica como OrcaBend.
%     \item Linguagens de programação utilizadas para implementação de modelos matemáticos como Python e MATLAB.
%     \item Métodos numéricos aplicados à análise de estruturas viscoelásticas, incluindo método de diferenças finitas (FDM), Runge-Kutta e método da perturbação.
%   \end{itemize}
%   \item Apresentação das séries de Prony como ferramenta fundamental para desenvolvimento da formulação matemática e implementação computacional da análise estrutural do \emph{bend stiffener}.
% \end{itemize}

% ----------------------------------------------------------------------
% Pedro
\chapter{Objetivo Geral}
\label{cha:objetivo-geral}

O objetivo geral deste trabalho é aplicar e validar uma metodologia numérica
para caracterização viscoelástica de poliuretanos e análise de enrijecedores à
flexão para \emph{risers} flexíveis, empregando modelos constitutivos identificados a 
partir de ensaios de relaxação e fluência (\emph{creep}), implementação computacional
da formulação de análise e proposição de critérios para dimensionamento de 
enrijecedores à flexão para um carregamento harmônico simples.

\section{Objetivos Específicos}
\label{sec:objetivos-especificos}

O trabalho tem como objetivos específicos:

\begin{enumerate}
  \item Realizar revisão bibliográfica sobre: viscoelasticidade, poliuretanos, 
enrijecedores à flexão, modelos físicos e séries de Prony.
  \item Implementar procedimentos para identificação de parâmetros de modelos 
viscoelásticos a partir de dados de ensaios de relaxação.
  \item Implementar numericamente a formulação viscoelástica para enrijecedores à flexão, 
  considerando os efeitos viscoelásticos do poliuretano e carregamento harmônico simples.
  \item Avaliar e validar a resposta estrutural do enrijecedor e comparar com os resultados 
  esperados pela bibliografia.
  \item Realizar um estudo de caso.
  \item Analisar os resultados obtidos no estudo de caso e discutir sobre as limitações 
  envolvendo a metodologia.

\end{enumerate}

\section{Premissas Adotadas} % Mover para metodologia

As premissas adotadas para a realização do modelo matemático de enrijecedor e tubo flexível
desenvolvido neste trabalho incluem: 
\begin{itemize}
  \item A Teoria de Euler-Bernoulli para vigas e a consideração de um engaste fixo na 
  extremidade do enrijecedor e da linha flexível.
  \item A deformação axial é desprezada, de modo que a linha flexível e o enrijecedor se deformam apenas por flexão.
  \item A linha flexível que reside no interior do enrijecedor é considerada como parte integrante do sistema, sem descontinuidade de deslocamentos.
  \item O comportamento em tração e compressão é considerado o mesmo, tanto para a linha flexível
  quanto para o enrijecedor.
  \item A força e o ângulo de topo são modelados como funções harmônicas devido à natureza do ambiente
  a qual o sistema está inserido, isto é, de carregamentos induzidos por ondas.
  \item O comportamento de tração e compressão são considerados os mesmos.
  \item O efeito da temperatura sobre as propriedades do poliuretano é desprezado.
\end{itemize}

Embora as simplificações negligenciem alguns dos efeitos termofísicos importantes, como a redução 
de capacidade (\emph{derating}) associado à variação de temperatura e a anisotropia do poliuretano,
essas simplificações permitem uma análise consistente do comportamento viscoelástico do sistema como descrito
por Caire \cite{caire}.

% ---------------------------------------------------------------------
% Alexandre
\chapter{Referencial Teórico}
\label{cha:referencial-teorico}

Nesse capítulo apresenta-se o referencial teórico que fundamenta o desenvolvimento do trabalho, 
abordando conceitos essenciais relacionados à viscoelasticidade, poliuretanos, enrijecedores à 
flexão e modelos constitutivos utilizados na análise estrutural de \emph{risers} flexíveis.
Ademais, explicitam-se os conceitos básicos sobre equações diferenciais, métodos numéricos e 
modelos constitutivos utilizados na solução destas. Por fim, são discutidos os conceitos de 
séries de Prony e sua aplicação na caracterização viscoelástica de materiais poliméricos, 
com ênfase em poliuretanos.

\section{Flexão}
\label{sec:flexao}

\section{Viscoelasticidade}
\label{sec:viscoelasticidade}

\subsection{Rigidez}
\label{sub:rigidez}

\section{Equações Diferenciais}
\label{sec:equacoes-diferenciais}

\section{Métodos Numéricos}
\label{sec:metodos-numericos}

\section{Séries de Prony}
\label{sec:series-de-prony}

% ----------------------------------------------------------------------
% Pedro
\chapter{Metodologia}
\label{sec:metodologia}

\chapter{Resultados e Discussão}
\label{sec:resultados-e-discussao}

% ----------------------------------------------------------------------

\chapter{Conclusão}

\section{Análises}

\section{Trabalhos Futuros}


%%% ====================================================================
%%% Início da parte pós-textual do documento.

\postextual


%%% Referências Bibliográfica
\bibliography{referencias}

%%% Início dos apêndices ---------------------------------------------

\apendices

%%% A linha abaixo impede que as seções, subseções, etc. dos apêndices e anexos
%%% sejam mostradas no Sumário.
\addtocontents{toc}{\protect\setcounter{tocdepth}{0}}

\partapendices*

% --------------------------------------------------------------------

\chapter{Demonstração do Teorema Loren Ipsum}

% --------------------------------------------------------------------
%%% Início dos anexos ------------------------------------------------
\anexos

\partanexos*

\chapter{Exemplo de lombada}

\hspace*{0mm}\fbox{\includegraphics[scale=0.8]{pics/normas-ifes-apendice-d.png}}

\end{document}